\begin{textsample}{POS Dim 1 – 2014 – Score 25.00 – 2014/t001657.txt}  \label{ex:f1_pos_045}
On June 12, 2014, precisely at 3:33 in a balmy winter afternoon in SĂŁo Paulo, Brazil, a typical South American winter afternoon, this kid, this young man that you see celebrating here like he had scored a goal, Juliano Pinto, 29 years old, accomplished a magnificent deed.
Despite being paralyzed and not having any sensation from mid-chest to \textbf{the} tip of his toes as \textbf{the} result of a car crash six years ago that killed his brother and produced a complete spinal cord lesion that left Juliano in a wheelchair, Juliano rose to \textbf{the} occasion, and on this day did something that pretty much everybody that saw him in \textbf{the} six years deemed impossible.
Juliano Pinto delivered \textbf{the} opening kick of \textbf{the} 2014 Brazilian World Soccer Cup here just by thinking.
He could not move his body, but he could imagine \textbf{the} movements needed to kick a ball.
He was an athlete before \textbf{the} lesion.
He ' s a para-athlete right now.
He ' s going to be in \textbf{the} Paralympic Games, I hope, in a couple years.
But what \textbf{the} spinal cord lesion did not rob from Juliano was his ability to dream.
And dream he did that afternoon, for a stadium of about 75, 000 people and an audience of close to a billion watching on TV.
And that kick crowned, basically, 30 years of basic research studying how \textbf{the} brain, how this amazing \textbf{universe} that we have between our ears that is only comparable to \textbf{universe} that we have above our head because it has about 100 billion elements talking to each other through \textbf{electrical} brainstorms, what Juliano accomplished took 30 years to imagine in laboratories and about 15 years to plan.
When John Chapin and I, 15 years ago, proposed in a paper that we would build something that we called a brain-machine interface, meaning connecting a brain to \textbf{devices} so that animals and humans could just move these \textbf{devices}, no matter how far they are from their own bodies, just by imagining what they want to do, our colleagues told us that we actually needed professional help, of \textbf{the} psychiatry variety.
And despite that, a Scot and a Brazilian persevered, because that ' s how we were raised in our respective countries, and for 12, 15 years, we made demonstration after demonstration suggesting that this was possible.
And a brain-machine interface is not rocket science, it ' s just brain research.
It ' s nothing but using sensors to read \textbf{the} \textbf{electrical} brainstorms that a brain is producing to generate \textbf{the} motor commands that have to be downloaded to \textbf{the} spinal cord, so we projected sensors that can read hundreds and now thousands of these brain cells simultaneously, and extract from these \textbf{electrical} signals \textbf{the} motor planning that \textbf{the} brain is generating to actually make us move into space.
And by doing that, we converted these signals into digital commands that any mechanical, electronic, or even a virtual \textbf{device} can understand so that \textbf{the} subject can imagine what he, she or it wants to make move, and \textbf{the} \textbf{device} obeys that brain command.
By sensorizing these \textbf{devices} with lots of different types of sensors, as you are going to see in a moment, we actually sent messages back to \textbf{the} brain to confirm that that voluntary motor will was being enacted, no matter where — next to \textbf{the} subject, next door, or across \textbf{the} \textbf{planet}.
And as this message gave feedback back to \textbf{the} brain, \textbf{the} brain realized its goal : to make us move.
So this is just one experiment that we published a few years ago, where a monkey, without moving its body, learned to control \textbf{the} movements of an avatar arm, a virtual arm that doesn ' t exist.
What you ' re listening to is \textbf{the} sound of \textbf{the} brain of this monkey as it explores three different visually \textbf{identical} \textbf{spheres} in virtual space.
And to get a reward, a drop of orange juice that monkeys love, this animal has to \textbf{detect}, select one of these \textbf{objects} by touching, not by seeing it, by touching it, because every time this virtual hand touches one of \textbf{the} \textbf{objects}, an \textbf{electrical} \textbf{pulse} goes back to \textbf{the} brain of \textbf{the} animal describing \textbf{the} fine texture of \textbf{the} \textbf{surface} of this \textbf{object}, so \textbf{the} animal can judge what is \textbf{the} correct \textbf{object} that he has to grab, and if he does that, he gets a reward without moving a muscle. \textbf{The} perfect Brazilian lunch : not moving a muscle and getting your orange juice.
So as we saw this happening, we actually came and proposed \textbf{the} idea that we had published 15 years ago.
We reenacted this paper.
We got it out of \textbf{the} drawers, and we proposed that perhaps we could get a human being that is paralyzed to actually use \textbf{the} brain-machine interface to regain mobility. \textbf{The} idea was that if you suffered — and that can happen to any one of us.
Let me tell you, it ' s very sudden.
It ' s a millisecond of a collision, a car accident that transforms your life completely.
If you have a complete lesion of \textbf{the} spinal cord, you cannot move because your brainstorms cannot reach your muscles.
However, your brainstorms continue to be generated in your head.
Paraplegic, quadriplegic patients dream about moving every night.
They have that inside their head. \textbf{The} problem is how to get that code out of it and make \textbf{the} movement be created again.
So what we proposed was, let ' s create a new body.
Let ' s create a robotic vest.
And that ' s exactly why Juliano could kick that ball just by thinking, because he was wearing \textbf{the} first brain-controlled robotic vest that can be used by paraplegic, quadriplegic patients to move and to regain feedback.
That was \textbf{the} original idea, 15 years ago.
What I ' m going to show you is how 156 people from 25 countries all over \textbf{the} five \textbf{continents} of this beautiful Earth, dropped their lives, dropped their \textbf{patents}, dropped their dogs, wives, kids, school, jobs, and congregated to come to Brazil for 18 months to actually get this done.
Because a couple years after Brazil was awarded \textbf{the} World Cup, we heard that \textbf{the} Brazilian government wanted to do something meaningful in \textbf{the} opening ceremony in \textbf{the} country that reinvented and perfected soccer until we met \textbf{the} Germans, of course.
But that ' s a different talk, and a different neuroscientist needs to talk about that.
But what Brazil wanted to do is to showcase a completely different country, a country that values science and technology, and can give a gift to millions, 25 million people around \textbf{the} world that cannot move any longer because of a spinal cord injury.
Well, we went to \textbf{the} Brazilian government and to FIFA and proposed, well, let ' s have \textbf{the} kickoff of \textbf{the} 2014 World Cup be given by a Brazilian paraplegic using a brain- controlled exoskeleton that allows him to kick \textbf{the} ball and to feel \textbf{the} contact of \textbf{the} ball.
They looked at us, thought that we were completely \textbf{nuts}, and said,"Okay, let ' s try."We had 18 months to do everything from zero, from scratch.
We had no exoskeleton, we had no patients, we had nothing done.
These people came all together and in 18 months, we got eight patients in a routine of training and basically built from nothing this guy, that we call Bra-Santos Dumont 1. \textbf{The} first brain-controlled exoskeleton to be built was named after \textbf{the} most \textbf{famous} Brazilian scientist ever, Alberto Santos Dumont, who, on October 19, 1901, created and flew himself \textbf{the} first controlled airship on air in Paris for a million people to see.
Sorry, my American friends, I live in North Carolina, but it was two years before \textbf{the} Wright Brothers flew on \textbf{the} coast of North Carolina.
Flight control is Brazilian.
So we went together with these guys and we basically put this exoskeleton together, 15 degrees of freedom, hydraulic machine that can be commanded by brain signals recorded by a non-invasive technology called electroencephalography that can basically allow \textbf{the} patient to imagine \textbf{the} movements and send his commands to \textbf{the} controls, \textbf{the} motors, and get it done.
This exoskeleton was covered with an \textbf{artificial} skin invented by Gordon Cheng, one of my greatest friends, in Munich, to allow sensation from \textbf{the} joints moving and \textbf{the} foot touching \textbf{the} ground to be delivered back to \textbf{the} patient through a vest, a shirt.
It is a smart shirt with micro-vibrating elements that basically delivers \textbf{the} feedback and fools \textbf{the} patient ' s brain by creating a sensation that it is not a machine that is carrying him, but it is he who is walking again.
So we got this going, and what you ' ll see here is \textbf{the} first time one of our patients, Bruno, actually walked.
And he takes a few seconds because we are setting everything, and you are going to see a \textbf{blue} light cutting in front of \textbf{the} helmet because Bruno is going to imagine \textbf{the} movement that needs to be performed, \textbf{the} \textbf{computer} is going to analyze it, Bruno is going to certify it, and when it is certified, \textbf{the} \textbf{device} starts moving under \textbf{the} command of Bruno ' s brain.
And he just got it right, and now he starts walking.
After nine years without being able to move, he is walking by himself.
And more than that — — more than just walking, he is feeling \textbf{the} ground, and if \textbf{the} speed of \textbf{the} exo goes up, he tells us that he is walking again on \textbf{the} sand of Santos, \textbf{the} beach resort where he used to go before he had \textbf{the} accident.
That ' s why \textbf{the} brain is creating a new sensation in Bruno ' s head.
So he walks, and at \textbf{the} end of \textbf{the} walk — I am running out of time already — he says,"You know, guys, I need to borrow this thing from you when I get married, because I wanted to walk to \textbf{the} priest and see my bride and actually be there by myself.
Of course, he will have it whenever he wants.
And this is what we wanted to show during \textbf{the} World Cup, and couldn ' t, because for some mysterious reason, FIFA cut its broadcast in half.
What you are going to see very quickly is Juliano Pinto in \textbf{the} exo doing \textbf{the} kick a few minutes before we went to \textbf{the} pitch and did \textbf{the} real thing in front of \textbf{the} entire crowd, and \textbf{the} lights you are going to see just describe \textbf{the} operation.
Basically, \textbf{the} \textbf{blue} lights pulsating indicate that \textbf{the} exo is ready to go.
It can receive thoughts and it can deliver feedback, and when Juliano makes \textbf{the} decision to kick \textbf{the} ball, you are going to see two streams of green and \textbf{yellow} light coming from \textbf{the} helmet and going to \textbf{the} legs, representing \textbf{the} mental commands that were taken by \textbf{the} exo to actually make that happen.
And in basically 13 seconds, Juliano actually did.
You can see \textbf{the} commands.
He gets ready, \textbf{the} ball is set, and he kicks.
And \textbf{the} most amazing thing is, 10 seconds after he did that, and looked at us on \textbf{the} pitch, he told us, celebrating as you saw,"I felt \textbf{the} ball."And that ' s priceless.
So where is this going to go?
I have two minutes to tell you that it ' s going to \textbf{the} limits of your imagination.
Brain-actuating technology is here.
This is \textbf{the} latest : We just published this a year ago, \textbf{the} first brain-to-brain interface that allows two animals to exchange mental messages so that one animal that sees something coming from \textbf{the} environment can send a mental SMS, a torpedo, a neurophysiological torpedo, to \textbf{the} second animal, and \textbf{the} second animal performs \textbf{the} act that he needed to perform without ever knowing what \textbf{the} environment was sending as a message, because \textbf{the} message came from \textbf{the} first animal ' s brain.
So this is \textbf{the} first demo.
I ' m going to be very quick because I want to show you \textbf{the} latest.
But what you see here is \textbf{the} first rat getting informed by a light that is going to show up on \textbf{the} \textbf{left} of \textbf{the} cage that he has to press \textbf{the} left cage to basically get a reward.
He goes there and does it.
And \textbf{the} same time, he is sending a mental message to \textbf{the} second rat that didn ' t see any light, and \textbf{the} second rat, in 70 percent of \textbf{the} times is going to press \textbf{the} left lever and get a reward without ever experiencing \textbf{the} light in \textbf{the} retina.
Well, we took this to a little higher limit by getting monkeys to collaborate mentally in a brain net, basically to donate their brain activity and combine them to move \textbf{the} virtual arm that I showed you before, and what you see here is \textbf{the} first time \textbf{the} two monkeys combine their brains, synchronize their brains perfectly to get this virtual arm to move.
One monkey is controlling \textbf{the} x \textbf{dimension}, \textbf{the} other monkey is controlling \textbf{the} y \textbf{dimension}.
But it gets a little more interesting when you get three monkeys in there and you ask one monkey to control x and y, \textbf{the} other monkey to control y and z, and \textbf{the} third one to control x and z, and you make them all play \textbf{the} game together, moving \textbf{the} arm in 3D into a target to get \textbf{the} \textbf{famous} Brazilian orange juice.
And they actually do. \textbf{The} black dot is \textbf{the} average of all these brains working in parallel, in real time.
That is \textbf{the} definition of a biological \textbf{computer}, interacting by brain activity and achieving a motor goal.
Where is this going?
We have no idea.
We ' re just scientists.
We are paid to be children, to basically go to \textbf{the} edge and discover what is out there.
But one thing I know : One day, in a few decades, when our grandchildren \textbf{surf} \textbf{the} Net just by thinking, or a mother donates her eyesight to an autistic kid who cannot see, or somebody speaks because of a brain-to-brain bypass, some of you will remember that it all started on a winter afternoon in a Brazilian soccer field with an impossible kick.
Thank you.
Thank you.
Bruno Giussani : Miguel, thank you for sticking to your time.
I actually would have given you a couple more minutes, because there are a couple of points we want to develop, and, of course, clearly it seems that we need connected brains to figure out where this is going.
So let ' s connect all this together.
So if I ' m understanding correctly, one of \textbf{the} monkeys is actually getting a signal and \textbf{the} other monkey is reacting to that signal just because \textbf{the} first one is receiving it and transmitting \textbf{the} neurological impulse.
Miguel Nicolelis : No, it ' s a little different.
No monkey knows of \textbf{the} existence of \textbf{the} other two monkeys.
They are getting a visual feedback in 2D, but \textbf{the} task they have to accomplish is 3D.
They have to move an arm in three \textbf{dimensions}.
But each monkey is only getting \textbf{the} two \textbf{dimensions} on \textbf{the} video screen that \textbf{the} monkey controls.
And to get that thing done, you need at least two monkeys to synchronize their brains, but \textbf{the} ideal is three.
So what we found out is that when one monkey starts slacking down, \textbf{the} other two monkeys enhance their performance to get \textbf{the} guy to come back, so this adjusts dynamically, but \textbf{the} global synchrony remains \textbf{the} same.
Now, if you flip without telling \textbf{the} monkey \textbf{the} \textbf{dimensions} that each brain has to control, like this guy is controlling x and y, but he should be controlling now y and z, instantaneously, that animal ' s brain forgets about \textbf{the} old \textbf{dimensions} and it starts concentrating on \textbf{the} new \textbf{dimensions}.
So what I need to say is that no Turing machine, no \textbf{computer} can predict what a brain net will do.
So we will absorb technology as part of us.
Technology will never absorb us.
It ' s simply impossible.
BG : How many times have you tested this?
And how many times have you succeeded versus failed?
MN : Oh, tens of times.
With \textbf{the} three monkeys?
Oh, several times.
I wouldn ' t be able to talk about this here unless I had done it a few times.
And I forgot to mention, because of time, that just three weeks ago, a European group just demonstrated \textbf{the} first man-to-man brain-to-brain connection.
BG : And how does that play?
MN : There was one bit of information — big ideas start in a humble way — but basically \textbf{the} brain activity of one subject was transmitted to a second \textbf{object}, all non-invasive technology.
So \textbf{the} first subject got a message, like our rats, a visual message, and transmitted it to \textbf{the} second subject. \textbf{The} second subject received a magnetic \textbf{pulse} in \textbf{the} visual cortex, or a different \textbf{pulse}, two different \textbf{pulses}.
In one \textbf{pulse}, \textbf{the} subject saw something.
On \textbf{the} other \textbf{pulse}, he saw something different.
And he was able to verbally indicate what was \textbf{the} message \textbf{the} first subject was sending through \textbf{the} Internet across \textbf{continents}.
Moderator : Wow. \textbf{Okay}, that ' s where we are going.
That ' s \textbf{the} next TED Talk at \textbf{the} next conference.
Miguel Nicolelis, thank you.
MN : Thank you, Bruno.
Thank you.

% matched lemmas: artificial, blue, computer, continent, detect, device, dimension, electrical, famous, identical, left, nut, object, okay, patent, planet, pulse, sphere, surf, surface, the, universe, yellow
\end{textsample}
